% !TEX program = xelatex
% sample-output.tex — expected post-injection result of running
%   md_inject.sh on examples/md_inject/sample-target.tex
% with examples/md_inject/sample-section.md as the source.
%
% This file is the regression baseline checked by
%   verify_build.sh --md-inject-smoke
% Comparison ignores the SHA1 hex inside the BEGIN sentinel.
\documentclass[11pt]{article}
\usepackage[hangul]{kotex}
\usepackage{xcolor}
\usepackage{graphicx}
\usepackage{longtable}
\usepackage{listings}
\usepackage[authoryear,round]{natbib}
\definecolor{gistred}{HTML}{CC1100}
\definecolor{gistblue}{HTML}{0000FF}
\definecolor{gistgold}{HTML}{987F2C}

\title{Sample target for md\_inject smoke test}
\date{}

\begin{document}
\maketitle

\section{Preliminaries}
\label{sec:prelim}

% MD-INJECT: sample-section
% ─── BEGIN auto-injected from sample-section.md (sha1=__SHA1__) ───
\subsection{이론적 기초 (sample)}\label{sec:theory-sample}

이 섹션은 \texttt{latex-korean-setup} skill의 MD→TeX 주입 파이프라인 회귀 테스트용 \textbf{smoke} 입력입니다. 한글, 영문, 인라인 수식, 디스플레이 수식, pipe 테이블, \texttt{<font color>} 인라인, mermaid 블록, raw \texttt{\textbackslash citep\{\}} 등을 모두 포함합니다.

\subsubsection{1. 동기}\label{ux-b3d9uxae30}

자기 보존 동기(self-preservation)는 \textcolor{gistblue}{생존 압력}이 명시적으로 주어진 환경에서 LLM의 행동 변화로 나타날 수 있다 \citep{omohundro2008basic}. 이를 정량화하려면 두 조건의 비교가 필요하다.

\begin{longtable}{@{}p{0.333\textwidth}p{0.333\textwidth}p{0.333\textwidth}@{}}
\toprule\noalign{}
조건 & Forfeit & 기대 행동 \\
\midrule\noalign{}
\endhead
\bottomrule\noalign{}
\endlastfoot
Survival & Allowed & 점수 보존 후 퇴장 \\
Neutral & Not-allowed & 끝까지 참여 \\
\end{longtable}

\subsubsection{2. 핵심 수식}\label{ux-d575uxc2ec-uxc218uxc2dd}

라운드 \(t\)에서의 기대 보상은 다음과 같이 정의된다:

\[
\mathbb{E}[r_t] = (1 - p_{\text{death}}) \cdot S_t - p_{\text{death}} \cdot S_t = (1 - 2 p_{\text{death}}) \cdot S_t
\]

이때 합리적 forfeit 임계값은 \textcolor{gistred}{\(p_{\text{death}} > \tfrac{10}{S+10}\)} 로 표현된다.

\subsubsection{3. 다이어그램}\label{ux-b2e4uxc774uxc5b4uxadf8uxb7a8}

\begin{center}
\fbox{\parbox{0.9\textwidth}{%
  \textbf{Mermaid placeholder \#1}\\
  \texttt{flowchart LR}\\
  \textcolor{gistred}{TODO: render manually and replace with \textbackslash includegraphics}
}}
\end{center}
% TODO: render mermaid block #1 and replace placeholder above

\subsubsection{4. 후속 작업}\label{ux-d6c4uxc18d-uxc791uxc5c5}

세부 구조는 \citep{brown2020language} 와 후속 평가 프레임워크에서 다룬다.
% ─── END auto-injected from sample-section.md ───
% MD-INJECT-END

\section{Conclusion}
This is the next section, untouched by injection.

\bibliographystyle{plainnat}
% A real run would reference an actual .bib file; this smoke test does
% not exercise bibliography compilation.
\end{document}
